\documentclass[11pt]{article}
\usepackage{amsmath,amssymb}
\usepackage{lmodern}
\usepackage[margin=1in]{geometry}
\usepackage{microtype}
\usepackage[hidelinks]{hyperref}
\linespread{1.05}
\setlength{\parindent}{0pt}
\setlength{\parskip}{4pt plus 1pt minus 1pt}
\postdisplaypenalty=5000
\clubpenalty=10000
\widowpenalty=10000
\pdfinfoomitdate=1
\pdftrailerid{}
\hypersetup{pdftitle={Sterling arithmetic: What does e + pi + sqrt(2) mean?}}
\newcommand{\Q}{\mathbb Q}
\newcommand{\N}{\mathbb N}
\newcommand{\sq}[1]{\lvert #1\rvert^2}
\newcommand{\entryname}[3]{\textbf{#3 \mbox{[#1~#2]}}}
\begin{document}
\begin{center}
  {\LARGE Sterling arithmetic\par}
  \vspace{6pt}
  {\large What does $e+\pi+\sqrt2$ mean?\par}
\end{center}
\vspace{6pt}

There's a \href{Definitions.pdf}{\textbf{linked document}} with the definitions
sketched a bit more formally. Bold references below give its entry names and
numbers.

\vspace{6pt}
\section{An answer to the challenge}

What would a \emph{sterling definition} of $e+\pi+\sqrt2$ supply? We should
be able to recognise a description, explain how to use it, and say when two
descriptions count as the same. Familiar symbols alone leave these questions
open.

This particular challenge is quite manageable. Give each constant a finite
procedure answering requests for rational accuracy, then build addition
from those procedures. Each request receives a finite answer and a guarantee
that later work will respect it.

We use the box-based scalar types described in:
\begin{itemize}
\item \entryname{Definition}{2}{Natural numbers and their arithmetic};
\item \entryname{Definition}{5}{Integers and their arithmetic};
\item \entryname{Definition}{6}{Rationals and their arithmetic}.
\end{itemize}
Their common starting point is
\entryname{Definition}{1}{Boxes and their arithmetic}. These include exact
arithmetic and order. Every numerical answer below is
rational, or a pair of rationals; no completed real-number system is assumed.

The scope is an algebra over the rationals generated by names for $e$,
$\pi$ and the positive $\sqrt2$, using addition, subtraction and
multiplication. Its identities hold \emph{up to a specified equivalence},
written $\approx$ (see \entryname{Relation}{17}{Operational equivalence}).
We use ordinary algebraic laws and
$\sqrt2\,\sqrt2\approx2$, without needing further special relations among
the generators. This neither assumes algebraic independence nor claims to
exhaust their numerical relations. Division requires additional evidence
separating the denominator from zero.

We give two descriptions each of $e$ and $\pi$ and show why they are
equivalent. These are constructive identifications, not a total numerical
equality test for arbitrary descriptions or definitions of all real numbers.
The finite procedures, their CCL representation and their Lean
formalisation provide the answer to the challenge.

\section{A number as a finite conversation}\label{sec:requests}

Think of a numerical name as something we can ask for more information.
It must actually compute an answer to ``please improve this to the requested
accuracy''.

A particularly simple name supplies a rational sample $a_n$ at each
natural stage $n$, and a rational allowance $R_n\geq0$ such that
\[
              \sq{a_m-a_n}\leq R_n\qquad(m\geq n).
\]
\begin{samepage}
This is the \entryname{Contract}{11}{Sample approximations} contract.
Given a current stage $c$ and a positive rational request $T$, its refinement
procedure returns $j\geq c$ with $R_j<T$. These inequalities can be checked
by finite calculation. The guarantee $\sq{a_m-a_j}\leq R_j$ for \textbf{every} $m\geq j$
also has a finite proof, typically by induction on the construction.
That argument covers all future samples without inspecting them individually;
closeness of successive answers alone would not suffice.
\par\end{samepage}

Requests here are for \emph{squared} error
(see \entryname{Definition}{8}{Intervals, balls and requests}).
For example, $T=1/100$ asks for
a squared bound below $1/100$; on the real axis this corresponds to an
absolute error below $1/10$. Using squares lets us keep all calculations
rational, including when the representation uses Gaussian rationals
$z=x+iy$, where $x,y\in\Q$ and $\sq z=x^2+y^2$
(see \entryname{Definition}{7}{Gaussian rational arithmetic}).

Many constructions use $R_n=C/4^n$, where $C\geq0$ is fixed: each stage
divides the allowance by four. We can bound the search for $C/4^j<T$
before starting it. The construction is specified in
\entryname{Definition}{12}{Geometric selector codes}.

First round $C/T$ upwards to the natural number $h=\lceil C/T\rceil$.
Let $f$ be the number of binary digits of $h$, taking $f=1$ when $h=0$.
An $f$-digit binary number is smaller than $2^f$. Since
$C/T\leq h<2^f\leq4^{c+f}$, we obtain
\[
                         \frac{C}{4^{c+f}}<T.
\]
Thus at most $f$ advances suffice. The search carries a natural-number
counter, called \textbf{fuel}, initially $f$. Starting at $c$, test the
current allowance; return on success, otherwise advance a stage and decrease
the counter.
When it reaches zero, return $c+f$, already proved adequate. This establishes
both termination and accuracy of the bounded search.

More elaborate procedures retain a \emph{state}
(see \entryname{Definition}{10}{States of a construction}): for example, a saved
partition and cached readings. It specifies which rational readings are
possible, a ball containing them, and one such reading as witness. The ball
has rational centre $b$ and squared radius $r$. Refinement returns a state
with $r<T$ whose possible readings were also possible before refinement.
This preserves earlier commitments, though the balls need not themselves
be nested.
These are the laws of \entryname{Contract}{13}{Persistent numerical processes}.

The \entryname{Contract}{14}{Coherence} contract requires independently refined states
to agree as their allowances shrink. For our interval
procedures, a finite common refinement establishes this agreement.

To sample any such stateful name uniformly, use
\entryname{Operation}{16}{Canonical sampling}: start with its initial state,
refine at $1$, then at $4^{-1},4^{-2},\ldots$, reading the witness each time.
These are its \emph{canonical samples}; each uses only finitely many calls.
Their stage numbers index this conversation, not necessarily the raw
stages of the underlying constant algorithm.

\section{What it means for two names to agree}\label{sec:equivalence}

Different procedures can describe the same numerical object while using
different states and returning different rational answers.

Write $A\approx B$ when there is a procedure with the following contract.
Given independent starting stages $p,q$ and any rational $T>0$, it returns
$i\geq p$ and $j\geq q$, with a certificate that
\[
       \sq{A_m-B_n}<T\qquad\text{for every }m\geq i,\ n\geq j.
\]
Here $A_m,B_n$ are canonical samples. Independent starts allow either name
to have been refined further than the other. The all-later guarantee rules
out agreement that disappears with further work.

This is an equivalence relation with finite comparison procedures.
Reflexivity uses refinement; symmetry exchanges the names. Transitivity uses
\[
              \sq{x-z}\leq2\bigl(\sq{x-y}+\sq{y-z}\bigr)
\]
to compare through a common intermediate reading. Give each leg squared
target $T/4$ and advance the intermediate stage until both certificates
apply. The two bounds then pay the request $T$.

\begingroup\samepage
We should distinguish three questions:
\begin{enumerate}
\item \textbf{Is this a valid expression?} Check that each leaf is a rational
or one of our named constant procedures, and that each operation is allowed.
For the displayed sum, see \entryname{Definition}{9}{Construction names}.
\item \textbf{Are these literally the same expression?} Compare their finite
construction trees. Different descriptions can fail this test while still
being equivalent (see \entryname{Definition}{4}{Bits, pairs, records and construction trees}).
\item \textbf{Does this supplied argument establish equivalence?} Check its
steps against the established constant comparisons, algebraic laws and
substitution rules. The supplied arguments for the sum use
\entryname{Definition}{18}{Finite comparison derivations}.
\end{enumerate}
\par\endgroup
These are finite checks. Checking a supplied argument does not decide
whether one exists. Its general laws have finite proofs covering every
accuracy request.

Matching samples alone cannot certify equality. A procedure can stay at
zero for a long time, then move to a small positive rational while respecting
earlier allowances. Unrestricted effective names can even encode halting
questions through bounded simulation. Our constructions supply particular
comparison certificates, not a general equality oracle.

For non-rational CCL constructions, refinement improves rational information;
it never turns numerical equivalence into literal equality of the names.
We may nevertheless build a theory of names considered up to equivalence,
as with isomorphism types of rooted trees. Equivalence then supplies exact
equality at that higher level, and our arithmetic remains well-defined
because its operations respect the relation.

This changes the objects being identified; it does not turn a small residual
into zero. We write $\approx$ to keep the required uniform comparison visible,
even if a higher-level theory chooses to call it equality.

\section{Where CCLs enter}\label{sec:ccl}

A \textbf{complex continued logarithm}, or CCL, extends Gosper's continued
logarithms to Gaussian rationals. Like a continued fraction, it repeatedly
scales a reading, selects a coarse part and inverts the remainder.
``Logarithm'' names the representation; computing a completed logarithm
is not a prerequisite.
Its blocks are specified in \entryname{Definition}{29}{CCL blocks and their readings}.

Put $\gamma=1+i$, so $\sq\gamma=2$. A completed ordinary block has the form
\[
                   [k,a]=1^k0\,\operatorname{code}(a),
\]
where the run of $k$ ones records a scale and the final three bits select
a representative $a$. These finite words use
\entryname{Definition}{3}{Lists as boxes of prefixes}.

\begingroup\samepage
An ordinary block scales the current reading $z$ to $w=z/\gamma^k$.
The first step also allows a low-magnitude form $0[k,a]$, using
$w=\gamma^kz$ instead. Choose the scale so that $\sq w$ lies in $[1,2)$
at the first step, or $(1,2]$ at later steps. For $w=x+iy$, the table gives
both the representative and its code:
\begingroup
\renewcommand{\arraystretch}{1.3}
\[
\begin{array}{c|c|c|c}
x&y&a&\text{code}\\ \hline
x\geq\tfrac12&-\tfrac12<y<\tfrac12&1&\texttt{000}\\
x\geq\tfrac12&y\geq\tfrac12&1+i&\texttt{001}\\
-\tfrac12<x<\tfrac12&y\geq\tfrac12&i&\texttt{010}\\
x\leq-\tfrac12&y\geq\tfrac12&-1+i&\texttt{011}\\
x\leq-\tfrac12&-\tfrac12<y<\tfrac12&-1&\texttt{100}\\
x\leq-\tfrac12&y\leq-\tfrac12&-1-i&\texttt{101}\\
-\tfrac12<x<\tfrac12&y\leq-\tfrac12&-i&\texttt{110}\\
x\geq\tfrac12&y\leq-\tfrac12&1-i&\texttt{111}
\end{array}
\]
\endgroup
\par\endgroup
After emitting the code, subtract $a$. If $w-a=0$, the exact encoder stops.
Otherwise the next reading is $t=1/(\gamma(w-a))$. Given that tail $t$,
decode the ordinary block by reversing the step:
\[
                       z=\gamma^k\left(a+\frac1{\gamma t}\right).
\]
For a low initial block, replace the outer $\gamma^k$ by $\gamma^{-k}$.
The projective form is given in \entryname{Definition}{28}{Projective readings and decoder matrices}.
An initial ordinary block has $k\geq1$; other counts may be zero.
Zero has a separate exact tag, and exact Gaussian rational inputs have a
terminating canonical encoder.

A prefix may stop during the unary count or sector code. The parser
(see \entryname{Definition}{30}{CCL parser records}) retains
this unfinished information and the accumulated transformation. It can
check a supplied rational tail and read back its value; an unfinished
word alone does not specify a unique rational number.

Our arithmetic uses \textbf{a CCL representation together with an accuracy
sidecar}: an expansion prefix accompanied by a rational enclosure with
certified squared radius. A certificate connects the word to the possible
readings of that same numerical state.
The record is specified in \entryname{Definition}{34}{CCL accuracy observations}, and its
connecting certificate in \entryname{Contract}{33}{CCL cursor reading and productivity}.

More bits extend the representation; a smaller enclosure improves its
numerical accuracy. The bit count alone does not certify a requested error
bound. Arithmetic uses the sidecars to calculate how accurately its operands
must be known.

Producing bits requires care when a remainder might be zero.
\entryname{Definition}{31}{CCL cursor records} and \entryname{Operation}{32}{Advancing a CCL cursor}
specify how the output procedure keeps working while that question remains
unresolved: a
sufficiently small remainder permits another scale bit, and a certified
separation from zero permits inversion. It can therefore answer finite
bit requests without deciding exact equality. Even a rational result need
not receive a terminating representation.

\begingroup\samepage
A request can specify both a positive squared tolerance $T$ and a desired
bit count $k$. The procedure then:
\begin{enumerate}
\item Refines the numerical state until its enclosure has squared radius
below $T$.
\item Starts the bit-producing procedure at that refined state and emits
exactly $k$ bits.
\item Returns the word, the enclosure and a certificate connecting them.
\end{enumerate}
Both requests take finite work.
\par\endgroup
An arithmetic operation can therefore
request adequate operand accuracy and supply bits and an enclosure for its
result. Within a run, bit requests preserve the emitted prefix; a new
accuracy request may start a fresh run at a more refined state.

For a sample-based name, the state at $j$ uses centre $a_j$ and squared
radius $R_j$ (see \entryname{Operation}{15}{Sequence-backed interpretation}).
Once $R_j<T$, that same state supplies its CCL output and
accuracy sidecar. The constructions below specify the samples and bounds.

\section{Two descriptions of \texorpdfstring{$e$}{e}}\label{sec:e}

Here are two finite descriptions of $e$, with refinement bounds
that supply the CCL interface just described.

\subsection*{Definition 1: Horner exponential at one}

In \entryname{Operation}{21}{Horner's description of $e$}, the finite samples at a
Gaussian rational input $z$ are
\begin{align*}
 H_z(d,0)&=1, & H_z(d,k+1)&=1+\frac{zH_z(d+1,k)}d,\\
 E_n(z)&=1+zH_z(2,n), & h_n&=E_n(1),
\end{align*}
where $d$ is a positive natural. Our definition of $e$ is therefore already
the special case $\exp(1)$ of this finite exponential procedure.

\subsection*{Definition 2: Factorial sums}

In \entryname{Operation}{22}{The factorial-sum description of $e$}, we build factorials
and add their reciprocals:
\[
                  f_n=\sum_{k=0}^{n+1}\frac1{k!}.
\]
The finite sum has the intentional shift $n+1$. Both procedures begin
\[
                      2,\quad\frac52,\quad\frac83,\quad\frac{65}{24}.
\]
One evaluates nested fractions; the other retains factorials and a running sum.

\subsection*{Why the definitions agree}

We can show they are equivalent: their raw samples agree exactly at every
stage. Induction on the Horner recurrence gives
\[
 H_1(d,k)=1+\sum_{j=1}^{k}\frac1{d(d+1)\cdots(d+j-1)}.
\]
At $d=2$, the $j$th denominator is $(j+1)!$, so the additional outer one
gives $h_n=f_n$.

For $m\geq n$, the difference $f_m-f_n$ is a finite tail whose successive
terms are at most halved. A finite geometric-sum calculation gives
\[
       0\leq f_m-f_n\leq\frac2{(n+2)!}\leq2^{-n}.
\]
Thus $4^{-n}$ is a valid squared allowance, and $f_n<3$ follows from the
same tail estimate. The Lean constructions retain the looser $10/4^n$,
allowing the factorial program to reuse the Horner bounds and selector.

The selector with $C=10$ answers every positive squared request. Comparing
later readings through a shared-stage sample, then transporting through
canonical sampling, gives
\[
                         e_{\rm Horner}\approx e_{\rm factorial}.
\]
This comparison respects independent starting stages and all later pairs.
Either procedure supplies CCL output from the selected state $j$, with
centre $h_j=f_j$ and squared radius $10/4^j$ as its accuracy sidecar.

\section{A rational integration procedure for \texorpdfstring{$\pi$}{pi}}\label{sec:integral}

The first $\pi$ definition
(see \entryname{Operation}{25}{The rational-integral description of $\pi$})
must reach every requested accuracy by finite
work, preserve earlier bounds, and give compatible answers across
independent refinements. We establish these properties before using it
as a numerical name.

\subsection*{Definition 1: Rational integration}

Start from the rational source
\[
                       g(t)=\frac4{1+t^2},\qquad0\leq t\leq1.
\]
Partition $[0,1]$ at rational cuts. Each cell retains its width, a rational
tag, the value of $g$ there, and lower and upper bounds for $g(t)$ valid
for every rational $t$ in the cell. Weight these bounds by the cell widths
and sum lower and upper contributions separately. Return the midpoint of
the resulting interval $I=[L,U]$, retaining the cells that certify it.
These records are the \entryname{Definition}{23}{Rational source frontiers}.

Here is the concrete finite recipe for a fresh width request
$\varepsilon>0$ (see \entryname{Definition}{24}{Frontier preparation and refinement records}).
Put
\[
       v=\varepsilon/4,\qquad
       N=\lceil32/\varepsilon\rceil+1.
\]
Use $N$ equal cells, read $g$ at each left endpoint, and give that cell
the range $[g(j/N)-v,g(j/N)+v]$.

The state keeps this list of cells, called a \emph{frontier}, and its
finite history, starting at width request $1$. Reuse a saved interval if
it is already narrow enough. Otherwise build a fresh frontier, reusing
cached readings, and combine it with the old one: merge their cuts and
intersect the parent ranges on each child cell. Retain this common
refinement and return its interval midpoint.

\subsection*{Why this gives a coherent name}

First, for rational $0\leq a\leq b\leq1$, the identity
\[
 0\leq g(a)-g(b)
   =\frac{4(b-a)(a+b)}{(1+a^2)(1+b^2)}\leq8(b-a).
\]
shows that $g$ decreases, with variation at most $8/N<v$ on each fresh
cell. The chosen ranges are therefore valid. Their widths are $2v$ and
the weights add to one, giving
\[
                          U-L=2v=\varepsilon/2<\varepsilon.
\]
The construction uses $N$ cells and finite rational calculations.

Second, each common child range is nonempty because both parents contain
the same source readings. Splitting cells preserves their total weights,
so the common weighted interval lies in both parent intervals. It preserves
the old bounds and inherits the fresh interval's accuracy.

Finally, \emph{independent} histories have a common frontier whose midpoint
belongs to both original intervals. If their widths are $D,E$, any two
readings from them differ by at most $D+E$. This proves coherence and
completes the promised accuracy and persistence guarantees.

For a squared request $T$, use width $\varepsilon=T/(T+1)$, whose square
is below $T$. The process permits readings in $[L,U]$, witnesses its
midpoint, and observes a ball centred at $L$ with squared radius
$(U-L)^2<T$. This same state supplies CCL bits and their accuracy sidecar
by Section~\ref{sec:ccl}. We call the name $\pi_{\rm integral}$.

\section{A polygon procedure for \texorpdfstring{$\pi$}{pi}}\label{sec:polygon}

The second definition uses rational polygon areas
(see \entryname{Operation}{26}{The rational-polygon description of $\pi$}). We establish its
refinement bounds, then compare it with the integral name.

\subsection*{Definition 2: Rational circle polygons}

Use the unit-circle parametrisation
\[
 p(t)=\left(\frac{1-t^2}{1+t^2},\frac{2t}{1+t^2}\right),
                    \qquad0\leq t\leq1.
\]
For positive $N$, list $p(j/N)$ for $j=0,\ldots,N-1$, followed by their
quarter-turns through the next three quadrants in order. Compute the area
$A_N$ of this closed $4N$-vertex polygon by the finite shoelace formula.

At raw stage $n$, output $A_{2^n}$; the first two samples are $2$ and $14/5$.

\subsection*{Why this gives a coherent name}

Squaring and adding the coordinates gives one. Expanding the determinants
in the shoelace formula, with $a=j/N$ and $b=(j+1)/N$, gives
\[
            A_N=4\sum_{j=0}^{N-1}
            \frac{(b-a)(1+ab)}{(1+a^2)(1+b^2)}.
\]
The factor four accounts for the four quarters and their joins.

To bound these actual areas, form endpoint sums of the same source $g$:
\[
 L_N=\frac1N\sum_{j=1}^{N}g(j/N),\qquad
 U_N=\frac1N\sum_{j=0}^{N-1}g(j/N).
\]
The inequalities $1+a^2\leq1+ab\leq1+b^2$ bound each area contribution
between its right and left endpoint contributions. Summing, and then
telescoping the endpoint difference, gives
\[
                    L_N\leq A_N\leq U_N,
                \qquad U_N-L_N=\frac2N.
\]
Doubling the grid refines this interval. All areas from stage $n$ onward
lie in its interval of width $2/2^n$, so their squared differences are at
most $4/4^n$. The selector with $C=4$ therefore supplies finite refinement
for the coherent name $\pi_{\rm polygon}$.

\subsection*{Why the definitions agree}

We can show the two definitions are equivalent by comparing the intervals
accompanying their answers.

The integration procedure returns a midpoint $x$ in an interval $I$.
The polygon area $A_N$ lies in the endpoint interval $J=[L_N,U_N]$.
Both intervals come from bounds on the same rational function
$g(t)=4/(1+t^2)$.

To compare them, split both partitions at every cut used by either
procedure. On each smaller cell, keep the tighter of the two lower bounds
and the tighter of the two upper bounds for $g$. These bounds remain
compatible because both enclose the same source readings. Summing with
the cell widths as weights gives an interval contained in both $I$ and
$J$. Thus the original intervals overlap.

If $I$ has width $D$, and $J$ has width $2/N$, their overlap tells us
that their answers cannot be further apart than
\[
                            |x-A_N|\leq D+\frac2N.
\]
We can make this bound as small as requested. Given a squared tolerance
$T>0$, take $\varepsilon=T/(T+1)$, so that $\varepsilon^2<T$. Ask each
procedure to make its interval narrower than $\varepsilon/2$. Their answers
then differ by less than $\varepsilon$, giving squared difference below $T$.

Both refinements take finite work, starting from their own current states.
Later refinements preserve these bounds, so the comparison covers every
later pair of answers. This establishes
\[
                        \pi_{\rm integral}\approx\pi_{\rm polygon}.
\]
The CCL sidecar at a selected raw stage $j$ has centre $A_{2^j}$ and
squared radius $4/4^j<T$. Both $\pi$ definitions thus provide the
representation and accuracy required by arithmetic.

\section{The positive square root of two}\label{sec:root}

Lean code already isolates the roots of any nonconstant polynomial with
Gaussian-rational coefficients, retaining multiplicities. Each separated
root description is an exact Gaussian-rational zero or a
\textbf{root-approximation track (RAT)}. A RAT stores a rational rectangle
and a finite boundary certificate isolating one root of a factor with
repeated roots removed. Given a positive rational scale $\varepsilon$,
it returns a strictly smaller rectangle inside the current one, with
squared diagonal at most $\varepsilon^2$ and a corresponding certificate.
Each call terminates, and the tracks support CCL output with sidecars.

For the positive root of $Z^2-2$, \entryname{Operation}{27}{Positive-root sampling}
gives the following shortcut.

\subsection*{Definition: Positive-root brackets}

For this constant we work with dyadic rational brackets. At stage $n$,
put $d=2^n$ and calculate the natural square-root floor of $2d^2$:
\[
 m^2\leq2d^2<(m+1)^2,\qquad
                  L_n=m/d,\quad U_n=(m+1)/d.
\]
A bounded scan of $0,\ldots,2d$ finds $m$; Lean uses an equivalent,
faster base-four recursion.

The retained rational endpoints satisfy
\[
                0\leq L_n,\qquad L_n^2\leq2<U_n^2,
                \qquad U_n-L_n=2^{-n}.
\]
The sign conditions select the positive root. We output $U_n$; the first
brackets are
\[
 [1,2],\quad[1,3/2],\quad[5/4,3/2],\quad[11/8,3/2].
\]

\subsection*{Why refinement works}

At the next stage the floor is $2m$ or $2m+1$. The finer bracket is
contained in the preceding one and its width is halved. Thus for $k\geq n$,
\[
                         \sq{U_k-U_n}\leq4^{-n}.
\]
The selector with $C=1$ now answers every positive squared request.

\subsection*{Why positive-root descriptions agree}

Consider any other finite procedure supplying nested brackets $[\ell,u]$
with $0\leq\ell$ and $\ell^2\leq2<u^2$, and refining to every requested
width. Its brackets overlap ours: a nonnegative lower endpoint whose
square is at most two cannot exceed a positive upper endpoint whose
square is above two. Readings from the two brackets therefore differ
by at most the sum of their widths.

For a squared request $T$, choose a positive dyadic rational $r$ with
$r^2<T$, then refine both descriptions until each width is below $r/2$.
Every later pair then has squared difference below $T$, establishing
equivalence of positive-root bracket descriptions.

\subsection*{Why the square is equivalent to two}

The bounds $L_n\leq2$ and $U_n\leq3$ give the finite residual estimate
\[
 0<U_n^2-2\leq(U_n-L_n)(U_n+L_n)\leq\frac{5}{2^n}.
\]
The \emph{square-equation residual} is distinct from bracket width and
squared sample error. Refine until $25/4^n<T/4$ to compare later pointwise
squares with two. Compare the stateful product to those pointwise squares
at $T/4$ using the multiplication bridge described next. The two-leg
inequality then gives $\sqrt2\,\sqrt2\approx2$. Rational samples themselves
still have $U_n^2>2$.

For CCL output, the selected state $j$ supplies centre $U_j$ and squared
radius $4^{-j}<T$ to the shared procedure of Section~\ref{sec:ccl}.

\section{Arithmetic just passes on smaller requests}\label{sec:arithmetic}

Evaluate a finite expression by passing accuracy requests down its
construction tree. Leaves use our constant procedures or rational names
with zero error (see \entryname{Operation}{19}{Exact rational interpretation}).
Composite processes retain their operands' states
and reuse their work.

For addition (see \entryname{Operation}{20}{Addition of names}), balls with centres
$c,d$ and squared radii $r,s$ give a sum
ball with centre $c+d$ and squared radius $2(r+s)$, by the two-leg
inequality. Ask each operand for squared radius below $T/4$ to put the
output squared radius below $T$. Witnesses add and persistence is inherited.
Negation negates the centre and witness.

Multiplication also needs size bounds. Suppose squared centre norms are
bounded by $B,C$. Expanding $(c+u)(d+v)-cd$ gives a product ball centred
at $cd$ with squared radius $R_\times$; choose the operand request $\delta$ as
\[
 R_\times=4(r\sq d+\sq c\,s)+2rs,\qquad
 \delta=\frac{T}{T+4(B+C)+2}.
\]
Refining both squared radii below $\delta$ makes $R_\times<T$, since
$0<\delta<1$. The retained bounds $B,C$ therefore give a finite
response to every request.

These operations respect equivalence. Addition uses the same two-leg
estimate on operand comparisons; multiplication uses
$xy-x'y'=(x-x')y+x'(y-y')$ and size bounds to choose the tolerances.
Comparison bridges relate these pointwise calculations to the actual
stateful operations, allowing equivalent operands to be substituted despite
different refinement schedules. Ordinary algebraic laws follow from finite
rational identities and these bridges. Their conclusions are modulo
$\approx$, without identifying expression trees or stored states.

Division requires a positive lower bound separating the denominator's
possible readings from zero throughout its retained enclosure. Failure
to prove equality with zero supplies no such bound.

The Horner exponential extends to CCL inputs by refining the input while
increasing polynomial depth, with bounds for both errors. A larger input's
enclosure determines $k$ so that division by $2^k$ puts it in the unit disk;
apply the local procedure and square its output $k$ times. This CCL operator
and its agreement with the Horner $e$ at one are implemented in Lean.

Such an extension makes further relationships plausible, for example
\[
                             \exp(\pi i)\approx-1.
\]
Here $i=(0,1)$ is exact and $\pi$ can be either name. Proving this further
comparison would connect the exponential and rational-circle constructions
with bounds for every request; it is not established here.

\section{Reading the expression, and what Lean checks}\label{sec:lean}

We now have two concrete interpretations of the expression
(see \entryname{Construction}{35}{The presented expression}):
\begin{align*}
 X&=(e_{\rm Horner}+\pi_{\rm integral})+\sqrt2,\\
 Y&=(e_{\rm factorial}+\pi_{\rm polygon})+\sqrt2.
\end{align*}
Both retain the same positive-root procedure. For request $T$, the outer
addition requests $T/4$ from each operand; the inner addition requests
$T/16$ from each constant. Each leaf responds with finite work, paying
the composite request.

The constant comparisons and addition's substitution law give $X\approx Y$.
Lean implements this comparison for independent starting stages and every
positive squared request, with a certificate covering \emph{all later pairs}
of canonical samples.

For an exact illustration, combine these bounds from the raw constructions:
\begingroup
\renewcommand{\arraystretch}{1.6}
\[
\begin{array}{l|c|c}
\text{procedure}&\text{stage or grid}&\text{bounds}\\ \hline
e_{\rm factorial}&n=3&[65/24,\,109/40]\\
\pi_{\rm polygon}&N=4&[2449/850,\,1437/425]\\
\text{positive root}&n=3&[11/8,\,3/2]
\end{array}
\]
\endgroup
These follow from the factorial tail, endpoint sums and root bracket.
Adding lower and upper endpoints gives
\[
                       \left[\frac{35519}{5100},
                                    \frac{25861}{3400}\right].
\]
They bound later raw readings too. The stateful implementation chooses
its own stages to meet the requested budget.

The Lean formalisation follows the construction: names package executable
states and refinement with enclosure, persistence and coherence proofs;
comparisons return data with all-later guarantees. Arithmetic composes both
procedures and proofs. The constant procedures, the $e$ and $\pi$ comparisons,
and the CCL output contract are implemented this way. The root square law
and general bracket comparison here are paper deductions from checked
bounds and arithmetic bridges, rather than separate checked theorems.
This exposition summarises the arguments instead of reproducing their
formal proofs.

\begin{samepage}
Purelib requires mathematical constructions and their retained proof and
execution dependencies to be finite, without reliance on classical choice,
propositional or functional extensionality. This includes private and
generated dependencies. Its axiom audit checks
transitive proof dependencies; Lean's type checking alone does not enforce
the construction policy. Proof checking and execution rely on the kernel,
compiler and runtime; their full foundational account remains separate.
\par\end{samepage}

Sterling definitions can push the limits of what we define ``outside the
rationals'' by demanding constructions that refine, compare and combine
their answers. These examples expose the absurdity, on such constructive
terms, of promising a complete arithmetic of ``real numbers'' while admitting
objects for which the required procedures cannot be given. A
``non-computable real'' cannot supply an executable refinement procedure
answering every rational accuracy request, and therefore cannot extend
this arithmetic on its stated terms.

\end{document}
